% !PRINTED 725
\weberpage{725}
%<<CONT
und folglich nach (10)
\[
N\varphi(\Theta) = 2^{\nu'-1}\,N\psi(\Theta).
\]
Daraus ergiebt sich
\weq{13}{A = N\psi(\Theta),}
und hieraus folgt mit Rücksicht auf (12), dass $A$ eine \gesperrt{ungerade
ganze Zahl ist.}

Bezeichnen wir die zu $m = 2^\varkappa$ gehörige Zahl $A$ mit $A_\varkappa$, so
ergiebt sich aus \S.\,197, (4) durch vollständige Induction der
Ausdruck für den ersten Factor der Classenzahl:
\weq{14}{k = A_4\,A_5\,\dots\,A_\varkappa,}
und daraus der Satz:

1. \gesperrt{Der erste Factor der Classenzahl ist eine ungerade ganze Zahl.}

Die Zahl $A_\varkappa$ ist nach der Formel (10) für die ersten Werthe
von $\varkappa$ nicht schwer zu berechnen.

Für $m = 16$ findet man $\nu' = 2$, $\beta = 0, 1$, $c_0 = c_1 = 1$
folglich
\[
\psi(\Theta) = 1, \quad A_4 = 1.
\]

Für $m = 32$, $\nu' = 4$, $\beta = 0, 1, 2, 3$.
\[
c_0 = 1, \quad c_1 = 1, \quad c_2 = -1, \quad c_3 = -1,
\]
\[
\psi(\Theta) = -\Theta^2,
\]
also auch $A_5 = 1$.

Für $m = 64$, $\nu' = 8$, $\beta = 0, 1, 2, 3, 4, 5, 6, 7$ sind die
absolut kleinsten Reste von $5^\beta$
\[
1,\ 5,\ 25,\ -3,\ -15,\ -11,\ 9,\ -19;
\]
also sind die $c_\beta$:
\[
+1, +1, +1, -1, -1, -1, +1, -1
\]
und
\[
\psi(\Theta) = -\Theta^3 + \Theta^6 - \Theta^7, \quad \psi(\Theta)\,\psi(-\Theta) = \Theta^6(\Theta^6 - 2i),
\]
woraus sich leicht ergiebt:
\[
A_6 = 17.
\]
Eine etwas längere Rechnung ergiebt noch
\[
A_7 = 21121.
\]
